% ===== §10 Conclusions: at the boundaries of the frameworks =====
\section{Conclusions: At the Boundaries of the Frameworks}
\label{sec:conclusions}

\noindent
A paper that argues that no theory possesses the whole of reason, that the truth of a relational being is generated in the dialogue among frameworks and settled by none, cannot in good faith end with a single, sovereign Conclusion. To do so would be to deny, in the performance, what it asserted in the proposition---a performative contradiction, and, in the paper's own vocabulary, a settling failure at the level of form: the use of one totalizing sign to declare the loop of the truth closed. This section therefore does not conclude. It offers, instead, \emph{Conclusions}---plural, local, mutually addressing, subsumed by no meta-framework---and closes on a reflection upon why it can offer no more.

A guard must be set at once, for the polyphonic gesture is easily corrupted into an evasion of the labour of argument---the false poem's ``staying open'' as an alibi (\xref{sec:false-poem}). True polyphony is not the refusal to assert; it requires that each voice sing to the full, to the very boundary of its own epistemic competence---each local conclusion hard, fallible, daring to assert---and then, honestly marking its limit, hand the microphone on. The test, here as throughout, is the holonomy: does the dialogue among the frameworks let the understanding climb in a spiral (each framework, illuminated by the others, carrying home a new remainder), or do mutually irrelevant assertions merely pile up in place? What follows tries to sing each voice to its boundary.

\subsection{The intermediate conclusion: a practical summation}
\label{sec:intermediate}

Before the frameworks each speak to their boundary, it is worth gathering, in one place, the practical result the paper has reached---an intermediate conclusion, offered not as the truth of the matter but as a summation of where the argument has arrived, so that the reader holds the whole before it is dispersed again into the polyphony.

For the sustainability of intimacy, the first thing is to \emph{create the poetic language}---the language that lets the generativity of value, across all three registers (the real, the symbolic, the imaginary), exist by sustaining itself rather than by being extracted. In this creation, the intervention of the big Other---of contract, of self-legislation (\xref{sec:contract})---is no less important: the manifold must be shaped, the floors set, the asymmetries legislated, so that the poetic flow has a just stage on which to run. As to how the poetic language is itself formed and applied: this too is something co-created and co-applied within a relational structure (\xref{sec:cocreation}), never presented ready-made by one party to another. And because the poetic language does not operate wholly within the symbolic register---reaching, as it must, into the real that no sign exhausts---we can give no mechanical, formalized theory of its perception or its practice (\xref{sec:no-metric}, \xref{sec:perception}). \emb{Indeed this may be among the central practices of intimacy itself: that in the very process of co-creation and interaction, the participants create the poetic language that lets the relation persist, and come to perceive the shared, unsymbolizable value that no formula could have delivered to them.} Upon the perception of that value, the poetic language built atop it lets the value circulate; and to let it circulate \emph{without extraction}---to let the value, in the very process of its generation, return to its creators themselves---is one way in which the justice of the cycle may be made to emerge. This, in sum, is the paper's practical result: not a method but a practice, not a structure presented but a language co-created, not a justice imposed but a justice that emerges when value is let to circulate and return. The frameworks will now say, each in its own tongue, what in this result it can and cannot underwrite.

\subsection{What the geometric framework can and cannot say}
\label{sec:concl-geom}

\emc{It can say:} that the difference between the good and the vicious cycle has an exact structural form---the circle of zero holonomy against the spiral of positive holonomy; that ``return to the root, yet not to the origin'' is no mere metaphor but the structure of anholonomy; that sublimation is the geometric phase accrued along a walked path, and therefore cannot be cashed out by a sign nor shortened by a shortcut; that depth is the area a path encloses. To this domain the framework brings a precision the others lack.

\emc{It cannot say:} how, from any external observation, to read off the sign of the curvature in a given case---the framework characterizes what the good circle \emph{is}, and offers no protocol for measuring it (\xref{sec:no-metric}); nor can it, of itself, fix the bearer of the phase (it had to borrow that from psychoanalysis), nor adjudicate the good (it had to hand that to justice). Its formalization is, by its own caveats (\xref{sec:caveats}), a structural analogy with a forward promise, not a completed physics. \emb{The geometric is one epistemic station among four, not the truth of the other three; formalization is not a higher truth, only a different tongue.} It hands the microphone to psychoanalysis for the bearer, and to justice for the good.

\subsection{What the psychoanalytic framework can and cannot say}
\label{sec:concl-psych}

\emc{It can say:} that the bearer of the phase is the sublimation of jouissance; that the good cycle encircles the void of \emph{das Ding} without filling it, and the vicious one fills it with an idol and so requires a sacrifice; that the sign anchors the real retroactively, as a \emph{point de capiton}; that real value, the least alienable, is bound up with the unsymbolizable and is approached only by encircling witness. It gives the paper its account of \emph{what is sublimated} and \emph{why the good cycle needs no enemy}.

\emc{It cannot say:} how its knowledge is itself secured, for that knowledge comes by way of transference, of \emph{Nachträglichkeit}, of the symptom---always partial, always resisted by the unconscious, never transparent to the subject who would wield it. It cannot deliver a self-present knowledge of one's own desire; the analyst's knowing is itself a knowing under erasure. \emb{Its truth is structurally partial, and it knows itself to be so.} It hands the microphone to political economy for the question of whose labour, materially, produces the phase that is here described in the register of desire.

\subsection{What the political-economic framework can and cannot say}
\label{sec:concl-poleco}

\emc{It can say:} that the good cycle returns value to all who create it and reduces no one to fuel; that the vicious cycle is an open system masquerading as a closed loop, sustaining a counterfeit phase by extraction from a hidden outside; that the alienation of social reproduction places the labour of care in the position of the invisible and uncounted; that the sign coordinates which equilibrium a system falls into, and that decentred reflux, not siphoning, is the mark of the good. It gives the paper its criterion of justice---whose blood the cycle runs on.

\emc{It cannot say:} that the agents whose incentives it analyzes are transparent to themselves, or that the structural account exhausts the meaning of what they do; the register of structural incentive is blind to the singular, unsubstitutable texture that the psychoanalytic and the poetic insist upon. It can tell us that a participant is being extracted from; it cannot, of itself, tell us what that participant's silence \emph{means} to her. \emb{It illuminates the structure and is blind to the singular.} It hands the microphone to the poem for what cannot be counted.

\subsection{What the Daoist framework can and cannot say}
\label{sec:concl-dao}

\emc{It can say:} that the default tendency of natural dynamics is toward the good, so that the burden of proof falls on intervention and not on letting-be; that the appropriate disposition is \emph{wu wei}, the guarding of a curvature that sublimates of itself; that the fear of an end, made into forced action, produces the end it dreads; that the good is, in a word, the dark virtue---to generate without possessing.

\emc{It cannot say:} this in the mode of a positive doctrine without betraying itself, for ``the one who knows does not speak'' (\zh{知者不言}); the moment its wisdom is objectified into a proposition to be asserted and defended, it has already falsified what it points to. Its truth is a knowing that withdraws from the saying---which is why it can ground the others' silences but cannot be made the master-discourse that subsumes them. \emb{It is the framework that knows the limit of all framing, and therefore the last that could claim to be the frame of frames.} It hands the microphone back to the silence from which the paper's saying must, in the end, return.

\subsection{Why this paper has no conclusion}
\label{sec:no-conclusion}

The four voices have sung, each to its boundary, and each has handed the microphone on rather than seized it. Note that none of the hand-offs returned to a fifth, higher voice that would gather them; the microphone passes among the four and to the silence, never up to a tribunal above them. This is not an incompleteness the paper regrets but the very thing it set out to enact.

\begin{claim}[The relationality of the conclusion]
A theory of relational being and sustainable generativity must have a conclusion that is itself relational. The four frameworks converge---on non-possession, on the spiral, on the return of value, on the dark virtue---yet none subsumes the others; the truth lives in the dialogue, the tension, the polyphony among them, and not in any one. To gather them into a single meta-framework's synthesis would be to perform upon the truth a settling failure---to cash out, with one sign, the understanding that can only be walked, again and again, in the frameworks' mutual rereading. There is no meta-language, no station outside all the frameworks from which to oversee them all---which is the methodological convergence of Lacan's ``there is no Other of the Other'' \citep{lacan2006} and the Daoist ``the one who knows does not speak.'' \emb{The paper therefore refuses a synthetic conclusion, and offers instead these plural, local, mutually inviting conclusions. The form of the paper is the demonstration of its content.}
\end{claim}

And so the paper ends where the prelude began, with the rose. The one who said the rose blooms but a week spoke of the flower; this paper has spoken of the root. The flower must fall---and must be let to fall, for the dread that will not let it fall is the dread that rots the root. What endures is not the flower held past its hour but the root that the fallen flower feeds: \emb{the non-possessive generativity, which perpetuates itself precisely because it does not clutch its own perpetuation.} This ninth paper has, in its turning, retroactively re-anchored the eighth---the lyric's address now revealed as the semiotic paradigm of the sustainable circle---which is itself a \emph{point de capiton}, a spiral return upon the series: the meaning settled by no single paper, generating new remainders along the path of the papers' mutual rereading. \zh{但愿人长久，千里共婵娟}: that we may endure, and share, though a thousand miles apart, the one fair moon---which no one possesses, and which returns its light to every upturned face.
